\chapter{Warnings and errors}\label{ch:trouble}

\begin{description}
\item[\texttt{steps of the deperturbed Dunham ladders are <= 0}] An
  excited level lies at or below a lower one. The log (\texttt{\#
  dunhamcheck:}) names the responsible constant and the cubic terms behind
  it.
  \begin{itemize}
  \item Type a (near Fermi resonance): use \opt{dunhamfix}~$=1$.
  \item Type b (a low-frequency, large-amplitude mode): use
    \opt{shiftgeometry}~$=1$ or \opt{samesignresonances}~$=1$.
  \end{itemize}
  See \notes{§11.8}.
\item[\texttt{rotatedcomp is ... > .5}] Some band mixes strongly with
  other states through terms treated as nonresonant, so its intensity is
  less accurate. The next line names the largest such term,
  with its harmonic detuning. If that detuning is small, make the term
  resonant: lower \opt{fermithres} (with \opt{rescriterion}~$=0$), or set its change to \texttt{on} in
  \file{resonancecandidates} and give the file back with
  \opt{resonanceoverrides}. With \opt{shiftgeometry} the rotation of the
  normal modes can also cause it.
\item[\texttt{freqsfile disables N modes in its 6th column}] That column
  was the choice of the fundamentals before 2026.11.1; it is now ignored.
  Choose the bands in \file{bandlist}.
\item[\texttt{error reading updown for band}] \file{dumpedparams} was
  written by the phase 1 of an older version. Rerun phase 1.
\item[\texttt{resintensitythres is no longer used}] Replace it by
  \opt{resintensityrel}.
\end{description}
